\documentclass[10pt,foldmarks=true,parskip=half]{scrlttr2}

\usepackage[ansinew]{inputenc}
\usepackage{multicol}
\usepackage{rotating}
\usepackage{xcolor}
\usepackage[german]{babel}
\usepackage[top=8mm, bottom=8mm]{geometry}
\usepackage{textcomp}
\usepackage{marvosym}
\usepackage{tabularx}
\usepackage{array}
\usepackage{booktabs}
\usepackage{relsize}

%%%herausgenommene Pakete
%\usepackage{fontenc}



%Paket f�r farbige Tabellen
\usepackage{colortbl}

%%%Serifenlose Schrift f�r das gesamte Dokument
\renewcommand{\familydefault}{\sfdefault}


%%Definition von Graut�nen
\definecolor{dunkelgrau}{rgb}{0.82,0.82,0.82}
\definecolor{hellgrau}{rgb}{0.92,0.92,0.92}

%Farbdefinitionen
\definecolor{dunkelgrau.80}{gray}{0.20}
\definecolor{dunkelgrau.60}{gray}{0.40}


%\KOMA-Optionen

%Schrift f�r Kopf- & Fusszeile anpassen:




%%%_________________________________________________________________________________________________

\begin{document}

\begin{letter}
 
%\pagestyle{empty}



	\begin{flushright}

		 


		\scriptsize

		\begin{tabular}{rl}
						 & \includegraphics[width=3cm]{DAK-Gesundheit_logo.png} \\
						 & \\
						 & \textbf{DAK} \\
						 & \textbf{22778 Hamburg} \\
		   \textbf{Telekontakt} 	 & Telefon: 030 911292980 \\
						 & Telefax: 030 9120298-7090\\
		   \textbf{Email}		 & service723300@dak.de \\
		   \textbf{Internet}		 & www.dak.de \\
		   \textbf{pers�nlicher Kontakt} & Tempelhofer Damm 158-160 \\
						 & 12099 Berlin \\


		\end{tabular} 

			     
		
	\end{flushright}


		\scriptsize

		  \setlength{\tabcolsep}{28pt}
		  \renewcommand{\arraystretch}{0.9}


		\noindent
		\begin{tabularx}{\textwidth}{@{}XXXX@{}}
		 \textbf{Datum} & \textbf{Unser Zeichen} 	& \textbf{Vorgang}	& \textbf{Sachbearbeiter(-in)} \\
		 02.04.2013	& 327 393 799 000-4		& 106560-13000-Pki	& Herr Hoffmann \\		\end{tabularx}

			      
\vspace{6mm}	

\normalsize



Pascal Bernhard \\
Schwalbacher Stra�e 7 \\
12161 Berlin

\vspace{18mm}


Sehr geehrter Herr Bernhard,\\

Sie hatten im Zeitraum vom 11. April 2011 bis einschlie�lich 30. September 2012 gem�� \S 44 Abs. 1  und \S 46 Abs. 2 SGB V von der DAK Krankengeld bezogen, da Sie aufgrund Ihrer Tumorerkrankung laut �rztlichem Gutachten vom 14. April 2011 als arbeitsunf�hig eingestuft wurden. Wir mussten nun feststellen, dass Sie zum Wintersemester 2011/2012 Ihr unterbrochenes Studium an der Freien Universit�t Berlin wieder aufgenommen haben und auch im darauffolgenden Sommersemester 2012 als regul�rer Student eingeschrieben waren. 

Eine Krankschreibung schlie�t jedoch ein Vollzeitstudium im gleichen Zeitraum aus nach \S 7 und \S 8 Abs. 1 SGB II. Da Sie ab dem 1. Oktober 2011 offenbar wieder studier- und somit auch arbeitsf�hig waren, lag ab diesem Datum die bescheinigte Arbeitsunf�higkeit nicht mehr vor. Entsprechend bestand im Zeitraum vom 1. Oktober 2011 bis 30. September 2012 kein Anspruch auf Krankengeld. Wir fordern Sie hiermit auf, die unberechtigt erhaltenen Leistungen innerhalb von 4 Wochen zur�ckzuerstatten. Den genauen Betrag entnehmen Sie bitte der folgenden Tabelle:



\vspace{8mm}


\small


%Tabelle
%\setlength{\tabcolsep}{15pt}
\renewcommand{\arraystretch}{1.4}


  \begin{tabularx}{\textwidth}{XXr}
    \rowcolor{dunkelgrau}

    \multicolumn{3}{l}{\textbf{Leistungen: Krankengeld} \scriptsize{\textsl{(nach V. Sozialgesetzbuch �44
Abs. 1 und �46 Abs. 2)}}} \\
    \hline
    \cellcolor{hellgrau} {\smaller{Bezugszeitraum:}} & \cellcolor{hellgrau}  {\smaller{monatlich}} & \cellcolor{hellgrau}  {\smaller{Summe}} \\

    {\smaller{\textsl{01.10.2011 -- 31.12.2011}}} & {\smaller{914,00}} \EUR & {\smaller{2745,00 \EUR{}}} \\
    {\smaller{\textsl{01.01.2012 -- 30.09.2012}}} & {\smaller{922,00 \EUR}} & {\smaller{8298,00 \EUR{}}} \\
    \hline
    {\smaller{\textbf{Gesamtbetrag:}}}		&		 & {\smaller{11043,00 \EUR{}}} \\
    
  \end{tabularx}




\par
\vspace{12mm}	


Mit freundlichen Gr��en,\\
\\
H. Hoffmann
			
Berlin, den 5. M�rz 2013

\vspace{2cm}



		\scriptsize

		  \setlength{\tabcolsep}{30pt}
		  \renewcommand{\arraystretch}{1.0}


		\noindent

		\begin{tabularx}{\textwidth}{@{}XXX@{}}
		 \toprule
		 \textcolor{dunkelgrau.60}{\textbf{Bankverbindung}} & \textcolor{dunkelgrau.60}{\textbf{Commerzbank Hamburg}} & \\
		 & \textcolor{dunkelgrau.60}{Konto: \emph{6202261}}	& \textcolor{dunkelgrau.60}{IBAN:  \emph{DE552004000006202261}}	\\
		  &  \textcolor{dunkelgrau.60}{BLZ: \emph{200 400 00}} & \textcolor{dunkelgrau.60}{BIC: \emph{COBA DE HHXXX}} \\

		\end{tabularx}
\newpage

	\begin{flushright}
	 
	\includegraphics[width=3cm]{DAK-Gesundheit_logo.png}

	\end{flushright}

\vspace{15mm}

\normalsize

\textbf{Rechtsbehelfsbelehrung:} \par

Gegen diesen Bescheid kann jeder Betroffene oder ein von diesem	bevollm�chtigter Dritter innerhalb eines Monats nach Bekanntgabe Widerspruch erheben. F�r Minderj�hrige oder nicht gesch�ftsf�hige Personen handelt deren gesetzlicher Vertreter. Der Widerspruch ist schriftlich oder zur Niederschrift bei der im Briefkopf genannten Stelle einzureichen.




\vspace{20cm}



		\scriptsize

		  \setlength{\tabcolsep}{30pt}
		  \renewcommand{\arraystretch}{1.0}


		\noindent

		\begin{tabularx}{\textwidth}{@{}XXX@{}}
		 \toprule
		 \textcolor{dunkelgrau.60}{\textbf{Bankverbindung}} & \textcolor{dunkelgrau.60}{\textbf{Commerzbank Hamburg}} & \\
		 & \textcolor{dunkelgrau.60}{Konto: \emph{6202261}}	& \textcolor{dunkelgrau.60}{IBAN:  \emph{DE552004000006202261}}	\\
		  &  \textcolor{dunkelgrau.60}{BLZ: \emph{200 400 00}} & \textcolor{dunkelgrau.60}{BIC: \emph{COBA DE HHXXX}} \\

		\end{tabularx}





\end{letter}


\end{document}

